% =====================================================================
%  Section 5: Discussion
% =====================================================================
\section{Discussion}
\label{sec:discussion}

\paragraph{What TADS delivers.}
The deliverable of \textsc{TADS} is the \emph{trajectory anchor}:
a model-intrinsic, per-layer capability direction extracted online
from the hidden-state trajectory of the very model being
fine-tuned. The anchor adds a structural signal to dynamic
re-ranking that scalar composite rewards alone cannot express, and
it is the object of both our method (it is what makes
\textsc{TADS} different from prior composite-reward selectors) and
our theory (Theorem~\ref{thm:stability} guarantees its stability
under a positive eigengap). Everything else in the selection rule
is built around the anchor: the composite reward
$\mathcal{R}_i^{(t)}$ (Eq.~\ref{eq:composite}) acts as a
stable carrier, and the single hyperparameter $\lambda$ controls
how strongly the anchor reshapes the ranking. Setting $\lambda=0$
disables the anchor and recovers the carrier alone, yielding a
clean ablation that isolates the anchor's contribution.

The anchor is model-intrinsic: it requires no seed examples,
validation labels, external judges, or auxiliary scorers, and it
is refreshed at every epoch as the model's hidden-state geometry
evolves. The stability guarantee we prove therefore concerns the
\emph{structural component} introduced by \textsc{TADS}, not the
composite-reward carrier or any particular scalar utility
transform layered on top.

\paragraph{Why the anchor helps.}
Dynamic re-ranking based only on a scalar transform of the
composite reward can be unstable in multi-task instruction pools.
In particular, the variance-ratio weight used to combine difficulty
and uncertainty drifts with task composition because it mixes
within-task and between-task variance
(Eq.~\ref{eq:total-variance}); the pool-level calibration removes
the additive baseline but does not address residual geometric
structure in the candidate distribution. The trajectory anchor adds
a pool-level geometric signal that favours samples aligned with the
representation direction the model is currently shaping.

This also separates \textsc{TADS} from dynamic data selection for
RLVR, such as DOTS~\cite{sun2025dots}. DOTS selects GRPO questions
using rollout outcomes and verifiable rewards; \textsc{TADS}
re-ranks supervised instruction--response pairs using hidden-state
geometry. DOTS is therefore an important related method, but not a
directly runnable baseline for our SFT setting, where automatic
verifiers, binary rewards, and rollout groups are unavailable by
default.

\paragraph{Implications and limitations.}
The stability bound makes the anchor a reproducible structural
object rather than probe noise. The eigengap $\gamma^{(t)}$ also
serves as a diagnostic: when it is large, the anchor is sharply
identified; when the covariance becomes nearly isotropic, the
anchor weakens and the method naturally reduces toward the
no-anchor composite-reward base.

\textsc{TADS} currently targets text-only SFT with single-sequence
instruction--response pairs and a rank-one anchor per layer. Future
extensions include turn-aware anchors for multi-turn dialogue,
separate anchors for multimodal representations, and rank-$r$
eigenspace anchors for samples combining multiple skills.
Extending trajectory anchoring to RL fine-tuning would require a
different experimental setup, since prompts, sampled responses,
rollouts, and reward-conditioned states define different possible
anchor spaces.
